% Options for packages loaded elsewhere
% Options for packages loaded elsewhere
\PassOptionsToPackage{unicode}{hyperref}
\PassOptionsToPackage{hyphens}{url}
\PassOptionsToPackage{dvipsnames,svgnames,x11names}{xcolor}
%
\documentclass[
  letterpaper,
  DIV=11,
  numbers=noendperiod]{scrartcl}
\usepackage{xcolor}
\usepackage{amsmath,amssymb}
\setcounter{secnumdepth}{-\maxdimen} % remove section numbering
\usepackage{iftex}
\ifPDFTeX
  \usepackage[T1]{fontenc}
  \usepackage[utf8]{inputenc}
  \usepackage{textcomp} % provide euro and other symbols
\else % if luatex or xetex
  \usepackage{unicode-math} % this also loads fontspec
  \defaultfontfeatures{Scale=MatchLowercase}
  \defaultfontfeatures[\rmfamily]{Ligatures=TeX,Scale=1}
\fi
\usepackage{lmodern}
\ifPDFTeX\else
  % xetex/luatex font selection
\fi
% Use upquote if available, for straight quotes in verbatim environments
\IfFileExists{upquote.sty}{\usepackage{upquote}}{}
\IfFileExists{microtype.sty}{% use microtype if available
  \usepackage[]{microtype}
  \UseMicrotypeSet[protrusion]{basicmath} % disable protrusion for tt fonts
}{}
\makeatletter
\@ifundefined{KOMAClassName}{% if non-KOMA class
  \IfFileExists{parskip.sty}{%
    \usepackage{parskip}
  }{% else
    \setlength{\parindent}{0pt}
    \setlength{\parskip}{6pt plus 2pt minus 1pt}}
}{% if KOMA class
  \KOMAoptions{parskip=half}}
\makeatother
% Make \paragraph and \subparagraph free-standing
\makeatletter
\ifx\paragraph\undefined\else
  \let\oldparagraph\paragraph
  \renewcommand{\paragraph}{
    \@ifstar
      \xxxParagraphStar
      \xxxParagraphNoStar
  }
  \newcommand{\xxxParagraphStar}[1]{\oldparagraph*{#1}\mbox{}}
  \newcommand{\xxxParagraphNoStar}[1]{\oldparagraph{#1}\mbox{}}
\fi
\ifx\subparagraph\undefined\else
  \let\oldsubparagraph\subparagraph
  \renewcommand{\subparagraph}{
    \@ifstar
      \xxxSubParagraphStar
      \xxxSubParagraphNoStar
  }
  \newcommand{\xxxSubParagraphStar}[1]{\oldsubparagraph*{#1}\mbox{}}
  \newcommand{\xxxSubParagraphNoStar}[1]{\oldsubparagraph{#1}\mbox{}}
\fi
\makeatother


\usepackage{longtable,booktabs,array}
\usepackage{calc} % for calculating minipage widths
% Correct order of tables after \paragraph or \subparagraph
\usepackage{etoolbox}
\makeatletter
\patchcmd\longtable{\par}{\if@noskipsec\mbox{}\fi\par}{}{}
\makeatother
% Allow footnotes in longtable head/foot
\IfFileExists{footnotehyper.sty}{\usepackage{footnotehyper}}{\usepackage{footnote}}
\makesavenoteenv{longtable}
\usepackage{graphicx}
\makeatletter
\newsavebox\pandoc@box
\newcommand*\pandocbounded[1]{% scales image to fit in text height/width
  \sbox\pandoc@box{#1}%
  \Gscale@div\@tempa{\textheight}{\dimexpr\ht\pandoc@box+\dp\pandoc@box\relax}%
  \Gscale@div\@tempb{\linewidth}{\wd\pandoc@box}%
  \ifdim\@tempb\p@<\@tempa\p@\let\@tempa\@tempb\fi% select the smaller of both
  \ifdim\@tempa\p@<\p@\scalebox{\@tempa}{\usebox\pandoc@box}%
  \else\usebox{\pandoc@box}%
  \fi%
}
% Set default figure placement to htbp
\def\fps@figure{htbp}
\makeatother


% definitions for citeproc citations
\NewDocumentCommand\citeproctext{}{}
\NewDocumentCommand\citeproc{mm}{%
  \begingroup\def\citeproctext{#2}\cite{#1}\endgroup}
\makeatletter
 % allow citations to break across lines
 \let\@cite@ofmt\@firstofone
 % avoid brackets around text for \cite:
 \def\@biblabel#1{}
 \def\@cite#1#2{{#1\if@tempswa , #2\fi}}
\makeatother
\newlength{\cslhangindent}
\setlength{\cslhangindent}{1.5em}
\newlength{\csllabelwidth}
\setlength{\csllabelwidth}{3em}
\newenvironment{CSLReferences}[2] % #1 hanging-indent, #2 entry-spacing
 {\begin{list}{}{%
  \setlength{\itemindent}{0pt}
  \setlength{\leftmargin}{0pt}
  \setlength{\parsep}{0pt}
  % turn on hanging indent if param 1 is 1
  \ifodd #1
   \setlength{\leftmargin}{\cslhangindent}
   \setlength{\itemindent}{-1\cslhangindent}
  \fi
  % set entry spacing
  \setlength{\itemsep}{#2\baselineskip}}}
 {\end{list}}
\usepackage{calc}
\newcommand{\CSLBlock}[1]{\hfill\break\parbox[t]{\linewidth}{\strut\ignorespaces#1\strut}}
\newcommand{\CSLLeftMargin}[1]{\parbox[t]{\csllabelwidth}{\strut#1\strut}}
\newcommand{\CSLRightInline}[1]{\parbox[t]{\linewidth - \csllabelwidth}{\strut#1\strut}}
\newcommand{\CSLIndent}[1]{\hspace{\cslhangindent}#1}



\setlength{\emergencystretch}{3em} % prevent overfull lines

\providecommand{\tightlist}{%
  \setlength{\itemsep}{0pt}\setlength{\parskip}{0pt}}



 


\usepackage{booktabs}
\usepackage{tabularx}
\usepackage{threeparttable}
\usepackage{ragged2e}
\KOMAoption{captions}{tableheading}
\makeatletter
\@ifpackageloaded{caption}{}{\usepackage{caption}}
\AtBeginDocument{%
\ifdefined\contentsname
  \renewcommand*\contentsname{Table of contents}
\else
  \newcommand\contentsname{Table of contents}
\fi
\ifdefined\listfigurename
  \renewcommand*\listfigurename{List of Figures}
\else
  \newcommand\listfigurename{List of Figures}
\fi
\ifdefined\listtablename
  \renewcommand*\listtablename{List of Tables}
\else
  \newcommand\listtablename{List of Tables}
\fi
\ifdefined\figurename
  \renewcommand*\figurename{Figure}
\else
  \newcommand\figurename{Figure}
\fi
\ifdefined\tablename
  \renewcommand*\tablename{Table}
\else
  \newcommand\tablename{Table}
\fi
}
\@ifpackageloaded{float}{}{\usepackage{float}}
\floatstyle{ruled}
\@ifundefined{c@chapter}{\newfloat{codelisting}{h}{lop}}{\newfloat{codelisting}{h}{lop}[chapter]}
\floatname{codelisting}{Listing}
\newcommand*\listoflistings{\listof{codelisting}{List of Listings}}
\makeatother
\makeatletter
\makeatother
\makeatletter
\@ifpackageloaded{caption}{}{\usepackage{caption}}
\@ifpackageloaded{subcaption}{}{\usepackage{subcaption}}
\makeatother
\usepackage{bookmark}
\IfFileExists{xurl.sty}{\usepackage{xurl}}{} % add URL line breaks if available
\urlstyle{same}
\hypersetup{
  pdftitle={Thesis Draft 2},
  colorlinks=true,
  linkcolor={blue},
  filecolor={Maroon},
  citecolor={Blue},
  urlcolor={Blue},
  pdfcreator={LaTeX via pandoc}}


\title{Thesis Draft 2}
\author{}
\date{}
\begin{document}
\maketitle


\section{Introduction}\label{introduction}

Over recent years, it has become increasingly clear that the scientific
community is facing a replication crisis (e.g., Forbes, Travers, and
Johnson 2023). A growing number of studies have reported difficulties
reproducing previously published findings. This problem is often
attributed to publication bias -- the tendency to only publish
significant results -- and questionable research practices (QRPs).
However, even studies aimed at removing the incentive for QRPs have
painted a rather pessimistic picture of reproducible science (Auspurg
and Brüderl 2024). Many-analyst studies, where multiple independent
researchers answer the same research question using the same dataset,
have shown considerable variation in analytic strategies and results
(e.g., Silberzahn et al. 2018; Bastiaansen et al. 2020; Breznau et al.
2022; Hoogeveen et al. 2023a; Gould et al. 2025; Kowall et al. 2025).
Reported effect sizes often vary from negative to positive, and from
significant to non-significant, all within the same study. This means
that there is not only considerable variation in the results but also in
the substantive conclusions drawn from them. At first glance, this seems
an immediate threat to the reproducibility of science. In an ideal
world, one might expect that a single dataset should yield a single,
`true' answer, raising the question whether effect size variation should
always be interpreted as problematic.

However, interpreting this variation without understanding its origins
may not be that straightforward. A frequently raised criticism of
many-analyst projects is that the posed research questions are often too
vague, leaving analysts uncertain about how to proceed with data
analysis and which variables should be included in their statistical
models (Kummerfeld and Jones 2023; Hoogeveen et al. 2023b).

\texttt{For\ example,\ @silberzahn\_many\_2018\ asked\ analysts\ to\ investigate\ the\ effect\ of\ skin\ tone\ of\ soccer\ players\ on\ the\ likelihood\ to\ receive\ a\ red\ card.\ This\ question\ can\ be\ interpreted\ different\ ways;\ are\ we\ interested\ in\ a\ descirptive\ association\ between\ skin\ tone\ and\ the\ likelihood\ of\ receiving\ a\ red\ card,\ or\ are\ we\ interested\ in\ a\ causal\ effect\ of\ skin\ tone\ on\ this\ likelihood?\ Depending\ on\ these\ different\ causal\ effects,\ different\ variables\ need\ to\ be\ included\ in\ the\ model.}

\emph{researchers may choose to identify different causal effects, each
of which has its merit, but each of which implies different estimands
and corresponding identification strategies (reasonable counterfactuals;
15, 36). For example, one might be interested in (i) identifying the
descriptive difference in the probability of receiving a red card by
skin color, or one might be interested in (ii) identifying the causal
direct effect of skin color (``racial bias'').}

This uncertainty is particularly problematic in situations where
researchers are implicitly interested in finding causal relationships,
which is often the case in many-analyst projects (Del Giudice and
Gangestad 2021; Bulbulia 2023; Rohrer, Smith, and Munafò 2025). In this
context, including different confounders may correspond to different
causal assumptions and thus different -- potentially opposing --
estimated effects. Therefore, when analysts select different variables,
they may in fact be answering different questions (Simonsohn, Simmons,
and Nelson 2015). From this perspective, effect size variation may
indicate meaningful substantive differences rather than problematic
analytical noise. On the other hand, if the effect size variation stems
from arbitrary, equally justifiable decisions whose options do not
reflect different substantive meanings, this may be considered more
problematic (Del Giudice and Gangestad 2021).

\subsubsection{The AB-framework}\label{the-ab-framework}

For better interpretation of effect size variation in many-analyst
projects, Auspurg and Brüderl (2024) have introduced a theoretical
framework (hereafter referred to as the AB-framework) to distinguish
between three different sources of variation, two of which are
considered `harmful' and one `helpful' (see
Table\textasciitilde{}\ref{tab:AB_framework}).

\begin{table}[H]
\centering
\caption{Sources of variation}
\label{tab:AB_framework}
\begin{threeparttable}
\footnotesize
\setlength{\tabcolsep}{10pt}
\renewcommand{\arraystretch}{1.2}

\begin{tabular}{
>{\RaggedRight\arraybackslash}p{7.2cm}
>{\RaggedRight\arraybackslash}p{7.2cm}
}
\toprule
\textbf{Harmful} & \textbf{Helpful} \\
\midrule

(a) Model uncertainty, including subjective analytical choices in data handling (e.g., missing data decisions, weighting etc.) and the selection of statistical models with different assumptions

\vspace{0.4em}
(b) Coding errors
&
(c) Substantively meaningful variation arising from different definitions of dependent, independent, and confounding variables, as well as differences in target populations (e.g., exclusion criteria or changes over time) \\

\bottomrule
\end{tabular}

\end{threeparttable}
\end{table}

First, effect size variation may arise from model uncertainty (a), for
example in missing-data handling or modelling assumptions, reflecting
ambiguity about the correct approach to estimate the parameter of
interest and its underlying data-generating process. For example, if
different missing data handling strategies (e.g., listwise deletion
versus multiple imputation) result in different mean effect sizes, while
the underlying observed data remain the same, this would imply that the
outcome depends on an arbitrary analytic choice rather than on the data
themselves. Secondly, effect size variation may be the result of
sloppiness or hidden coding errors (b) made by the analysts. Both
sources can be considered harmful, as they can introduce undesirable
uncertainty into scientific analyses, making conclusions less credible.

In contrast, the third source concerns analytical choices that provide
information about substantive variation that actually exists in the
population (c). As mentioned earlier, selecting different target
populations or confounders can lead to slightly different questions
being answered. For example, if analyst X uses the full dataset, while
analyst Y analyses only the men in that same dataset, X might find no
effect in the general population, while Y finds that there is a
male-specific positive effect. Given that the two analysts answered
different questions, this source can be considered as helpful for the
advancement of theories. By taking a closer look at the decisions made
by analysts and the implications thereof, the AB-framework provides
tools to interpret effect size variation in many-analyst projects.

\subsubsection{Identifying sources of variation in many-analyst
projects}\label{identifying-sources-of-variation-in-many-analyst-projects}

Although several many-analyst studies have aimed to identify the sources
of effect size variation, so far no consensus has been reached (e.g.,
Silberzahn et al. 2018; Bastiaansen et al. 2020; Breznau et al. 2022;
Hoogeveen et al. 2023a; Gould et al. 2025; Kowall et al. 2025). For
example, Hoogeveen et al. (2023a) found that the operationalisations of
variables, rather than model choice, influenced the results the most. In
contrast, Gould et al. (2025) showed that neither variation in variable
selection nor model structure meaningfully predicted variability, but
did not explicitly name another explanation source of variation. A
similar pattern was found by Kowall et al. (2025), where data structure
(cross-sectional versus longitudinal), rather than model or variable
choice, likely had the largest impact on the results. Different studies
thus highlight different sources of variation.

Given the wide variety of topics, designs, and variables across
projects, it may be unrealistic to expect a single, fixed set of
explanatory factors that applies to all studies. However, even
identifying a study-specific source of variation might not be possible
based on many-analyst studies alone. A key limitation of these projects
is that the involved analysts make decisions in an uncontrolled way
(Kummerfeld and Jones 2023). This leads to the problem of sparse data,
where little data is generally available for each individual decision,
as unique combinations of choices are often chosen only once. This
sparsity makes it difficult to properly disentangle individual sources
of variation and provide systematic quantitative evidence of their
relative influence. Consequently, solely looking at analytical decisions
in many-analyst projects may not provide enough information to interpret
effect size variation with the AB-framework.

\subsubsection{Multiverse analysis}\label{multiverse-analysis}

One approach to addressing the problem of data sparseness is multiverse
analysis (e.g., Steegen et al. 2016; Hoogeveen et al. 2023b). First,
this method identifies all reasonable decisions, each with a specific
number of options, and systematically combines them into a decision
space (Liu, Althoff, and Heer 2023). Since not all combinations may
yield possible analyses, the set of valid analysis paths (universes) is
always smaller than or equal to the full product. Then, results are
reported across all valid universes. This way, multiverse analysis
allows for the exploration of the entire decision space, rather than
just a subset of choices made by analysts. As a result, substantially
more data are available for each decision, enabling quantification of
their specific impact on the results.

Although multiverse analysis seems a promising way of systematically
evaluating the sensitivity of decisions, one of its limitations is that
the decision space is often defined by a single researcher (Simonsohn,
Simmons, and Nelson 2015). Consequently, some reasonable analysis paths
may be overlooked. However, by basing the multiverse analysis on an
existing many-analyst project, this problem can be solved. The decision
space will then be determined by the '\,'wisdom of the crowd'' and
represent realistic research choices (Auspurg and Brüderl 2024).

Another limitation is that different universes with contradicting
substantial implications may be included within the same multiverse. Del
Giudice and Gangestad (2021) argue that as a solution, a multiverse
should only include truly arbitrary, but not theoretically meaningful
decisions, as the latter may lead to misleading results. However, by
including both types of decisions and subsequently classifying them with
the AB-framework, results that might otherwise be considered misleading
could instead be interpreted in a meaningful way.

\subsubsection{Present study}\label{present-study}

To assess the current state of the replication crisis using many-analyst
projects, it is highly relevant to identify the most prominent sources
of variation and their implications based on the AB-framework. Due to
the descriptive nature of many-analyst studies, no consensus has yet
been reached about what factors drive variability in the results,
leading to the research question:

\emph{What is the relative contribution of harmful versus helpful
analytical decisions in many-analyst projects to effect size variation
when they are systematically evaluated in a multiverse?}

Aiming to gain quantitative evidence of the relative influence of
different sources of variation, the present study will conduct a
multiverse analysis based on the many-analyst project by Kowall et al.
(2025), where 16 teams analysed whether marital status (cohabiting
married versus never married) predicts cardiovascular disease. This was
chosen as a case study, because there was a large variety of both
harmful and helpful decisions, providing a rich decision space to
evaluate the impact of a wide range of different decisions. The decision
space consists of six analytical decisions: handling missing data,
choice of exclusion criteria, statistical model and definitions of
dependent variables, confounders and target population. Before running
the analysis, these decisions are classified as harmful/helpful using
the AB-framework. The multiverse analysis is largely based on sampling
and approximation methods proposed by Liu, Althoff, and Heer (2023),
where the decision space is sampled to immediately calculate the mean
effect and decision sensitivity. Lastly, a specification curve will be
plotted to visualise how the outcomes relate to different analysis
paths. For more details, see the Methods section.

\section{Methods}\label{methods}

\subsubsection{Ethics}\label{ethics}

Ethical approval was granted by the Ethics Review Board of the Faculty
of Social and Behavioural Sciences of Utrecht University.

\subsubsection{Data}\label{data}

Although the original teams' results are not analysed directly, the
multiverse is applied to the same data they analysed in Kowall et al.
(2025), namely the Survey of Health, Ageing and Retirement in Europe
(SHARE). This is a freely available (after registration) panel study
conducted in 28 European countries. Consistent with the original study,
waves 1-7 are used. As some analysts used the imputed SHARE datasets,
these are also used here.

All variables needed for the construction of the decision space (e.g.,
dependent, independent, confounders, exclusion) are selected from
different relevant datasets that are available per wave. Then, one
cleaned dataset is created per wave, where the selected variables are
linked through an ID variable.

\begin{table}[htbp]
\centering
\caption{Analytical decisions and their options based on the orginal many-analyst study, that were used in the present multiverse analysis.}
\label{tab:decision_space}
\begin{threeparttable}
\footnotesize
\setlength{\tabcolsep}{6pt}
\renewcommand{\arraystretch}{1.25}
\begin{tabularx}{\textwidth}{
>{\RaggedRight\arraybackslash}p{3.2cm}
>{\RaggedRight\arraybackslash}p{6.2cm}
>{\RaggedRight\arraybackslash}X
}
\toprule
\textbf{Decision} & \textbf{Options} & \textbf{Notes / Deviations from original} \\
\midrule
\multicolumn{2}{l}{\textbf{Harmful sources of variation}} \\
\midrule
Missing data &
Complete cases (16) \newline
``No information'' (1)
&
For complete cases, listwise deletion is applied. In the original study, analysts occasionally used the imputed dataset provided by SHARE, this was not done in this multiverse analysis. One team replaced missing values in categorical variables by ``no information'' to prevent listwise deletion. \\
\addlinespace
Model type &
Cox proportional hazards model (6) \newline
Logistic regression (5) \newline
Log-binomial model (2) \newline
Mixed-effects logistic regression (1) \newline
Discrete-time mixed-effects model (1) \newline
Generalised estimating equations (1)
&
Originally, different effect sizes were reported. For comparability, this multiverse analysis reports risk ratios only. \\
\addlinespace
\multicolumn{3}{l}{\textbf{Helpful sources of variation}} \\
\midrule
Data structure &
Strictly longitudinal (10) \newline
Strictly or partly cross-sectional (6) \newline
Longitudinal, wave 1 as only baseline (3) 
&
The cleaned datasets for each wave are combined in different ways, depending on the selected option. \\
\addlinespace
Exclusion criteria &
Subjects married before 18 \newline
Subjects from Israel \newline
Subjects under 55 at recruitment \newline
Subjects born after 1954
&
The analysts in the original study often used idiosyncratic exclusion criteria, which were then used in this study. Because the original paper did not specify whether combinations of these were used, a maximum of one exclusion criterion per universe was applied. \\
\addlinespace
Dependent variables &
Heart attack / stroke ever (16) \newline
Age at heart attack / stroke (9) \newline
Heart attack / stroke since last interview (7) \newline
Main cause of death: heart attack / stroke (5)
&
Seven combinations of the four dependent variables were used by the analysts in the original study. Although other combinations are possible, this decision only included the operationalisations from the original study. \\
\addlinespace
Confounders &
Age (12) \newline
Sex (12) \newline
Country (7) \newline
Smoking (6) \newline
Education (4) \newline
Physical activity (4) \newline
No confounders (2)
&
Only confounders used by at least four analysts, except `no confounders', were included to ensure feasibility of the multiverse analysis. \\
\bottomrule
\end{tabularx}
\begin{tablenotes}[flushleft]
\footnotesize
\item \textit{Note}. Numbers in parentheses indicate how often each option was selected across analytical teams.
\item \textit{Note}. For the independent variable (IV), all teams but two used ``cohabiting married" (CM). Some analysts also included registered partnership, and ``married living apart". However, since the original research question only focused on CM, this is the only variable used as IV.
\end{tablenotes}
\end{threeparttable}
\end{table}

\subsubsection{Definition of decision
space}\label{definition-of-decision-space}

To define the decision space, all options of six different decisions
made by the teams in Kowall et al. (2025) were identified based on
information provided in the original publication (see
Table\textasciitilde{}\ref{tab:decision_space}). When systematically
combining all options across decisions, the decision space consisted of
129024 universes. After removing invalid research paths, 82944 valid
universes remained.

\subsubsection{Classifying sources of variation using the
AB-framework}\label{classifying-sources-of-variation-using-the-ab-framework}

Before running the analyses, the selected decisions are classified using
the AB-framework. Decisions regarding missing data and statistical model
choice are considered harmful. Decisions regarding the
operationalisation of the dependent variable and the selection of
confounders and exclusion criteria are considered helpful. Additionally,
because the framework classifies variations over time as `different
target populations', which are meaningful, the data structure is
classified as helpful. Table \textasciitilde{}\ref{tab:decision_space}
also includes the classification of decisions.

\begin{figure}[H]
    \centering
    \includegraphics[width=1\linewidth]{round_robin_schematic.png}
    \caption{Schematic overview of Round Robin sampling algorithm. U: universe; T: the set of selected universes. Figure was created by the author for this research report.}
    \label{fig:round_robin}
\end{figure}

\subsubsection{Round Robin sampling}\label{round-robin-sampling}

Because the number of reasonable universes grows exponentially with the
number of decisions and their options, the computational cost of
multiverse analyses increases rapidly. To solve this problem, Liu,
Althoff, and Heer (2023) have proposed sampling universes one by one. By
directly calculating their outcomes and decision sensitivities, it is
possible to approximate multiverse results without having to explore the
full decision space. They evaluated various algorithms and found that
Round Robin sampling was one of the most effective methods. Therefore,
this method will be applied in the present study.

\textbf{NOTE: Because I eventually decided to program the Round Robin
sampler myself, I decided to move away from `calculating the outcomes
one by one'. I will include some text here describing how I programmed
the sampler in 2 steps: first by mapping what decision options
correspond to which row numbers, then sampling from these row numbers,
and subsequently filtering sampled row numbers from the original
decision space df (as sampling directly from the decision space was too
computationally heavy). This means that all universes are sampled
simultaneously and I ended up with a new file containing the sampled
decision space - which will be used for the multiverse. By setting a
seed, I ensure the sample is the same every time I run the sampler.}

Round Robin sampling aims to represent each decision option in the
multiverse at least once. The algorithm iteratively goes over all
decisions and their options and samples one universe at each stage. The
sampling is done without replacement, to make sure each sampled universe
is included at most once. For a visual example of how the algorithm
works, see Figure\textasciitilde{}\ref{fig:round_robin}.

When defining the number of sampled universes, it is important to make
sure enough data is available for every option to quantify the decision
sensitivity. The operationalisations of the dependent variable and the
confounders have the most options (7 and 64, respectively). Assuming a
sample size of at least 30 outcomes/universes per option (or more if
that is feasible computationally and within the time scope of this
project), this would result in 6400 sampled universes in total.

\subsubsection{Converting OR and HR to
RR}\label{converting-or-and-hr-to-rr}

\textbf{NOTE: In the first place, I wanted to convert all effect sizes
to Odds Ratio's. This turns out not to be possible. Based on:}

\begin{itemize}
    \item \url{https://pubmed.ncbi.nlm.nih.gov/31808145/}
    \item \url{https://ebm-helper.cn/en/Conv/OR_RR.html}
    \item \url{https://ebm-helper.cn/en/Conv/HR_RR.html}
\end{itemize}

\textbf{it seems that it is possible to convert both HR and OR to risk
ratios, which I will do. I will add a paragraph here explaining how this
conversion works.}

\subsubsection{Decision sensitivity}\label{decision-sensitivity}

So far, the robustness of multiverse outcomes has often been assessed
using specification curve analysis (SCA). In this type of analysis, each
universe is plotted against its associated estimated effect size
(Simonsohn, Simmons, and Nelson 2015). However, as this only provides
information about the robustness of the results, but not the decision
sensitivity, it will only be used as a visual aid to inspect if the full
outcome distribution is robust to different analytic strategies.

More recently, Liu, Althoff, and Heer (2023) have proposed using the
K-samples Anderson-Darling test (AD-test) to quantify not only the
robustness of the results, but also the decision sensitivity -- the
extent to which different options within each analytical decision lead
to different outcomes. For each option, an effect size distribution is
constructed, consisting of the results from all universes containing
that option. This means that within each outcome distribution the option
of one specific decision remains constant, while other analytical
decisions and their options vary. Next, the AD test is used to test
whether the outcome distributions of the different options within the
same decision originate from the same underlying distribution. If this
null hypothesis is rejected, the decision in question is classified as
`sensitive'.

Given that there are 64 confounder options, consisting of combinations
of 7 variables, it is uncertain whether the AD test will perform well.
Should this not be the case, the confounders will be analysed in a
slightly different way, namely by categorizing for each option whether
or not it is included in the universe, resulting in two outcome
distributions per option (included/excluded). In the case of the
confounders, this would mean that there would be only 14 options instead
of 64. A Kolmogorov-Smirnov (K-S) test, similar to the AD test but with
two samples, will then be used to quantify whether the two outcome
distributions originate from the same underlying distribution.

\subsubsection{Technical specifications}\label{technical-specifications}

All data are analysed using R, with packages \texttt{tidyverse},
\texttt{ggalluvial}, \texttt{multiverse}, \texttt{lme4},
\texttt{survival}, \texttt{geepack}, \texttt{kSamples} and \texttt{dgof}
Team (2024).

\subsubsection{Supplementary Materials}\label{supplementary-materials}

All code scripts containig sampler and multiverse functions can be found
\href{https://github.com/lodewijn/ManyAnalystThesisCompendium}{here}.

\section{Results}\label{results}

Figure\textasciitilde{}\ref{fig:sim_DV_dist} shows simulated example
distributions of outcomes from the universes across options of
\texttt{dependent\ variable\ operationalisations\textquotesingle{},\ 1-5\ were\ generated\ from\ a\ normal\ distribution,\ and\ 6-7\ from\ a\ beta-distribution.\ Given\ that\ not\ all\ outcome\ distributions\ have\ the\ same\ underlying\ distribution,\ the\ AD-test\ will\ \ classified\ this\ decision\ as}sensitive'.
Because this is a substantially meaningful decision, this would mean
that finding variation in the effect sizes may not be that problematic
for the reproducibility of science.

\section{Discussion}\label{discussion}

\section{Notes}\label{notes}

\begin{itemize}
\tightlist
\item
  liu: In other words, K--S test and LR may consider a decision to be
  more sensitive simply because the decision has more options.
\end{itemize}

\phantomsection\label{refs}
\begin{CSLReferences}{1}{0}
\bibitem[\citeproctext]{ref-arnold_dgof_2011}
Arnold, Taylor B., and John W. Emerson. 2011. {``Dgof: {Discrete}
{Goodness}-of-{Fit} {Tests}.''}
\url{https://CRAN.R-project.org/package=dgof}.

\bibitem[\citeproctext]{ref-auspurg_toward_2024}
Auspurg, Katrin, and Josef Brüderl. 2024. {``Toward a More Credible
Assessment of the Credibility of Science by Many-Analyst Studies.''}
\emph{Proceedings of the National Academy of Sciences} 121 (38):
e2404035121. \url{https://doi.org/10.1073/pnas.2404035121}.

\bibitem[\citeproctext]{ref-bastiaansen_time_2020}
Bastiaansen, Jojanneke A., Yoram K. Kunkels, Frank J. Blaauw, Steven M.
Boker, Eva Ceulemans, Meng Chen, Sy-Miin Chow, et al. 2020. {``Time to
Get Personal? {The} Impact of Researchers Choices on the Selection of
Treatment Targets Using the Experience Sampling Methodology.''}
\emph{Journal of Psychosomatic Research} 137 (October): 110211.
\url{https://doi.org/10.1016/j.jpsychores.2020.110211}.

\bibitem[\citeproctext]{ref-bates_fitting_2015}
Bates, Douglas, Martin Mächler, Ben Bolker, and Steve Walker. 2015.
{``Fitting {Linear} {Mixed}-{Effects} {Models} {Using} \textbf{Lme4}.''}
\emph{Journal of Statistical Software} 67 (1).
\url{https://doi.org/10.18637/jss.v067.i01}.

\bibitem[\citeproctext]{ref-breznau_observing_2022}
Breznau, Nate, Eike Mark Rinke, Alexander Wuttke, Hung H. V. Nguyen,
Muna Adem, Jule Adriaans, Amalia Alvarez-Benjumea, et al. 2022.
{``Observing Many Researchers Using the Same Data and Hypothesis Reveals
a Hidden Universe of Uncertainty.''} \emph{Proceedings of the National
Academy of Sciences} 119 (44): e2203150119.
\url{https://doi.org/10.1073/pnas.2203150119}.

\bibitem[\citeproctext]{ref-brunson_ggalluvial_2020}
Brunson, Jason. 2020. {``Ggalluvial: {Layered} {Grammar} for {Alluvial}
{Plots}.''} \emph{Journal of Open Source Software} 5 (49): 2017.
\url{https://doi.org/10.21105/joss.02017}.

\bibitem[\citeproctext]{ref-bulbulia_workflow_2023}
Bulbulia, Joseph A. 2023. {``A Workflow for Causal Inference in
Cross-Cultural Psychology.''} \emph{Religion, Brain \& Behavior} 13 (3):
291--306. \url{https://doi.org/10.1080/2153599X.2022.2070245}.

\bibitem[\citeproctext]{ref-del_giudice_travelers_2021}
Del Giudice, Marco, and Steven W. Gangestad. 2021. {``A {Traveler}'s
{Guide} to the {Multiverse}: {Promises}, {Pitfalls}, and a {Framework}
for the {Evaluation} of {Analytic} {Decisions}.''} \emph{Advances in
Methods and Practices in Psychological Science} 4 (1): 2515245920954925.
\url{https://doi.org/10.1177/2515245920954925}.

\bibitem[\citeproctext]{ref-forbes_supporting_2023}
Forbes, Heather J., Jason C. Travers, and Jenee Vickers Johnson. 2023.
{``Supporting the Replication of Your Research.''} In \emph{Research
{Ethics} in {Behavior} {Analysis}}, 237--62. Academic Press.
\url{https://doi.org/10.1016/B978-0-323-90969-3.00003-7}.

\bibitem[\citeproctext]{ref-gould_same_2025}
Gould, Elliot, Hannah S. Fraser, Timothy H. Parker, Shinichi Nakagawa,
Simon C. Griffith, Peter A. Vesk, Fiona Fidler, et al. 2025. {``Same
Data, Different Analysts: Variation in Effect Sizes Due to Analytical
Decisions in Ecology and Evolutionary Biology.''} \emph{BMC Biology} 23
(1): 35. \url{https://doi.org/10.1186/s12915-024-02101-x}.

\bibitem[\citeproctext]{ref-halekoh_r_2006}
Halekoh, Ulrich, Søren Højsgaard, and Jun Yan. 2006. {``The \emph{r}
{Package} \textbf{Geepack} for {Generalized} {Estimating}
{Equations}.''} \emph{Journal of Statistical Software} 15 (2).
\url{https://doi.org/10.18637/jss.v015.i02}.

\bibitem[\citeproctext]{ref-hoogeveen_many-analysts_2023}
Hoogeveen, Suzanne, Alexandra Sarafoglou, Balazs Aczel, Yonathan Aditya,
Alexandra J. Alayan, Peter J. Allen, Sacha Altay, et al. 2023a. {``A
Many-Analysts Approach to the Relation Between Religiosity and
Well-Being.''} \emph{Religion, Brain \& Behavior} 13 (3): 237--83.
\url{https://doi.org/10.1080/2153599X.2022.2070255}.

\bibitem[\citeproctext]{ref-hoogeveen_many-analysts_2023-1}
Hoogeveen, Suzanne, Alexandra Sarafoglou, Michiel van Elk, and Eric-Jan
Wagenmakers. 2023b. {``Many-Analysts Religion Project: Reflection and
Conclusion.''} \emph{Religion, Brain \& Behavior} 13 (3): 356--63.
\url{https://doi.org/10.1080/2153599X.2022.2070263}.

\bibitem[\citeproctext]{ref-kowall_marital_2025}
Kowall, Bernd, Linda Juel Ahrenfeldt, Jale Basten, Heiko Becher, Tilman
Brand, Julia Braun, Swaantje Casjens, et al. 2025. {``Marital Status and
Risk of Cardiovascular Disease -- a Multi-Analyst Study in
Epidemiology.''} \emph{European Journal of Epidemiology} 40 (5):
497--509. \url{https://doi.org/10.1007/s10654-025-01235-8}.

\bibitem[\citeproctext]{ref-kummerfeld_one_2023}
Kummerfeld, Erich, and Galin L. Jones. 2023. {``One Data Set, Many
Analysts: {Implications} for Practicing Scientists.''} \emph{Frontiers
in Psychology} 14 (February): 1094150.
\url{https://doi.org/10.3389/fpsyg.2023.1094150}.

\bibitem[\citeproctext]{ref-liu_approximation_2023}
Liu, Yang, Tim Althoff, and Jeffrey Heer. 2023. {``Approximation and
{Progressive} {Display} of {Multiverse} {Analyses}.''} arXiv.
\url{https://doi.org/10.48550/arXiv.2305.08323}.

\bibitem[\citeproctext]{ref-masur_specr_2020}
Masur, Philipp K., and Michael Scharkow. 2020. {``Specr: {Conducting}
and {Visualizing} {Specification} {Curve} {Analyses}.''}
\url{https://doi.org/10.32614/CRAN.package.specr}.

\bibitem[\citeproctext]{ref-rohrer_what_2025}
Rohrer, Julia M., George Davey Smith, and Marcus Munafò. 2025. {``What
Can Be Learned When Multiple Analysts Arrive at Different Estimates.''}
\emph{European Journal of Epidemiology} 40 (5): 493--95.
\url{https://doi.org/10.1007/s10654-025-01249-2}.

\bibitem[\citeproctext]{ref-sarma_multiverse_2021}
Sarma, Abhraneel, Alex Kale, Michael Jongho Moon, Nathan Taback, Fanny
Chevalier, Jessica Hullman, and Matthew Kay. 2021. {``Multiverse:
{Multiplexing} {Alternative} {Data} {Analyses} in {R} {Notebooks}.''}
Open Science Framework. \url{https://doi.org/10.31219/osf.io/yfbwm}.

\bibitem[\citeproctext]{ref-scholz_ksamples_2019}
Scholz, Fritz W., and Jiajuan Zhu. 2019. {``{kSamples}: {K}-{Sample}
{Rank} {Tests} and {Their} {Combinations}.''}
\url{https://CRAN.R-project.org/package=kSamples}.

\bibitem[\citeproctext]{ref-silberzahn_many_2018}
Silberzahn, R., E. L. Uhlmann, D. P. Martin, P. Anselmi, F. Aust, E.
Awtrey, Š. Bahník, et al. 2018. {``Many {Analysts}, {One} {Data} {Set}:
{Making} {Transparent} {How} {Variations} in {Analytic} {Choices}
{Affect} {Results}.''} \emph{Advances in Methods and Practices in
Psychological Science} 1 (3): 337--56.
\url{https://doi.org/10.1177/2515245917747646}.

\bibitem[\citeproctext]{ref-simonsohn_specification_2015}
Simonsohn, Uri, Joseph P. Simmons, and Leif D. Nelson. 2015.
{``Specification {Curve}: {Descriptive} and {Inferential} {Statistics}
on {All} {Reasonable} {Specifications}.''} \emph{SSRN Electronic
Journal}. \url{https://doi.org/10.2139/ssrn.2694998}.

\bibitem[\citeproctext]{ref-steegen_increasing_2016}
Steegen, Sara, Francis Tuerlinckx, Andrew Gelman, and Wolf Vanpaemel.
2016. {``Increasing {Transparency} {Through} a {Multiverse}
{Analysis}.''} \emph{Perspectives on Psychological Science} 11 (5):
702--12. \url{https://doi.org/10.1177/1745691616658637}.

\bibitem[\citeproctext]{ref-r_core_team_r_2024}
Team, R Core. 2024. {``R: {A} {Language} and {Environment} for
{Statistical} {Computing}.''} Vienna, Austria: R Foundation for
Statistical Computing. \url{https://www.R-project.org/}.

\bibitem[\citeproctext]{ref-therneau__a_2024}
Therneau, T. 2024. {``\_A {Package} for {Survival} {Analysis} in {R}\_.
{R} Package.''} \url{https://CRAN.R-project.org/package=survival}.

\bibitem[\citeproctext]{ref-wickham_welcome_2019}
Wickham, Hadley, Mara Averick, Jennifer Bryan, Winston Chang, Lucy
McGowan, Romain François, Garrett Grolemund, et al. 2019. {``Welcome to
the {Tidyverse}.''} \emph{Journal of Open Source Software} 4 (43): 1686.
\url{https://doi.org/10.21105/joss.01686}.

\end{CSLReferences}




\end{document}
